Las mezclas coherentes son aquellas que estan correlacionadas/son dependientes en tiempo y frecuencia,fase y estructura. Por ejemplo , dos guitarras tocando el mismo \textit{riff} en fase, o un coro cantando al unísono.

Las mezclas incoherentes son aquellas que tienen baja correlación, por ejemplo, una batería y una progresión de acordes que no concuerdan en el tono, por lo que en el espectro de frecuencias, ocupan rangos distintos.

Estas últimas son mas fáciles de separar por lo general ( depende de modelo y solapamiento ) pues los algoritmos pueden aprovechar esta diferencia frecuencial.